\section*{Bab \ref{ch:b2:wave-interfaces} --- Pemantulan dan Penerusan di Antarmuka}

\begin{solution}{exo:b2:wave-interfaces:1}
$r = (Z_2 - Z_1)/(Z_2 + Z_1)$, $R = r^2$: udara--air $0.9995$, $99.9\%$;
air--baja $0.93$, $86\%$; jaringan--tulang $0.63$, $40\%$; jaringan--udara
$-0.9995$, $99.9\%$. Suara Anda, di udara, dipantulkan permukaannya;
mesinnya mengguncang lambungnya, yang terbuat dari baja yang bersentuhan dengan airnya lalu
melewatkan $14\%$ getarannya lurus ke dalam.
\end{solution}

\begin{solution}{exo:b2:wave-interfaces:2}
$R = ((n - 1)/(n + 1))^2$: $2.0\%$, $4.3\%$, $17\%$; kaca--air $0.44\%$.
$0.957^{10} = 0.64$: sepertiga cahayanya hilang, sebagian besar sebagai bayangan hantu;
berlapis: $0.995^{10} = 0.95$.
\end{solution}

\begin{solution}{exo:b2:wave-interfaces:3}
$\tan\theta_B = n_2/n_1$: \qty{56.7}{\degree}, \qty{53.1}{\degree}, \qty{48.8}{\degree}. Cahaya
yang terpantul terpolarisasi tegak lurus bidang datangnya;
cahaya yang diteruskan hanya sebagian terpolarisasi (komponen s-nya
hanya sebagian terpantul). Pandanglah etalasenya pada \qty{57}{\degree} dari
normalnya dengan sumbu polarisatornya dalam bidang datangnya (mendatar
bagi etalase tegak yang dilihat dari samping).
\end{solution}

\begin{solution}{exo:b2:wave-interfaces:4}
Sudut kritisnya $\arcsin(1/n)$: \qty{41.1}{\degree}, \qty{48.8}{\degree}, \qty{24.4}{\degree}; $45^\circ > 41^\circ$:
prismanya memantulkan seluruhnya. $\ell = \lambda/2\pi\sqrt{n^2\sin^2\theta - 1}$: \qty{270}{nm}
pada \qty{45}{\degree}, \qty{115}{nm} pada \qty{60}{\degree}. Sudut kritis intan yang kecil
memerangkap cahaya yang memasukinya sampai ia keluar lewat faset atasnya ---
yaitu kecemerlangannya.
\end{solution}

\begin{solution}{exo:b2:wave-interfaces:5}
(a) $\sqrt{1.52} = 1.23$; $e = 550/4 \times 1.38 = \qty{100}{nm}$. (b) $r_1 = -0.160$, $r_2 =
-0.048$, $\varphi = \pi$: $r = (r_1 - r_2)/(1 - r_1r_2) = -0.112$, $R = 1.3\%$ (melawan
$4.3\%$). (c) \qty{400}{nm}: $\varphi = 1.375\pi$, $R \approx 2.2\%$; \qty{700}{nm}: $\varphi =
0.79\pi$, $R \approx 1.6\%$: biru dan merah lebih banyak terpantul daripada hijau ---
yaitu kilau keunguan sebuah lensa berlapis. (d) Beberapa lapisan meratakan $R$
di seluruh cahaya tampak.
\end{solution}

\begin{solution}{exo:b2:wave-interfaces:6}
(a) $T = 4 \times 30 \times 1.6/31.6^2 = 19\%$; sisanya berdengung, pulsanya
berdering. (b) $Z = \sqrt{Z_1Z_2} = \qty{6.9e6}{kg.m^{-2}.s^{-1}}$, $e = v/4f = \qty{0.15}{mm}$.
(c) $T_1 = 4 \times 30 \times 10^6 \times 410/(30 \times 10^6)^2 = 5.5 \times 10^{-5}$, $T_2 = 4 \times 1.6 \times
10^6 \times 410/(1.6 \times 10^6)^2 = 1.0 \times 10^{-3}$: $5.6 \times 10^{-8}$. (d) $4 \times 1.5 \times 1.6/
3.1^2 = 0.999$.
\end{solution}

\begin{solution}{exo:b2:wave-interfaces:7}
(a) $\omega_p/\omega = 3$: $\underline n = -\iu\sqrt8 = -2.83\iu$; $r = (1 + 2.83\iu)/(1 - 2.83
\iu)$, $|r| = 1$, fasenya $2\arctan2.83 = \qty{141}{\degree}$. (b) $c/\sqrt{\omega_p^2 - \omega^2}
= \qty{5.6}{m}$. (c) $E \propto \cos(kx + \varphi/2)$: simpul pertamanya di $kx = -(\pi + \varphi)/2$,
$x \approx \qty{-45}{m}$ di bawah lapisannya. (d) Indeksnya turun dengan ketinggian, jadi
Descartes membengkokkan sinarnya menjauhi arah tegak; ia berbalik di tempat
$n(z) = \sin\theta_0$, yakni $\omega_p^2 = \omega^2\cos^2\theta_0$, di bawah aras $\omega = \omega_p$;
frekuensi yang $\omega\cos\theta_0$-nya melampaui $\omega_p$ maksimum lapisannya
lolos.
\end{solution}

\begin{solution}{exo:b2:wave-interfaces:8}
$\sqrt{8\varepsilon_0\omega/\gamma}$: $2.7 \times 10^{-4}$ ($R = 0.9997$), $1.5 \times 10^{-2}$ ($0.985$),
$0.06$ ($0.94$) --- tetapi pada \qty{5e14}{Hz} $\omega\tau = 80$: tembaganya menjadi \omterm{def:b2:waves-in-media:plasma}{plasma},
bukan penghantar ohmik, lalu memantulkan sekitar $60\%$, berwarna. \qty{15}{W}
terserap dari berkas CO$_2$-nya: beberapa ratus kelvin tiap menit pada
cermin kecil, dari situlah airnya.
\end{solution}

\begin{solution}{exo:b2:wave-interfaces:9}
$r = 0.2$, $-1/3$, $-1$, $+1$, $0$; SWR-nya $1.5$, $2$, $\infty$, $\infty$, $1$; yang terpantul
$4\%$, $11\%$, $100\%$, $100\%$, $0$. \qty{96}{W} mencapai antenanya, \qty{4}{W}
kembali ke pemancarnya (dan sebagian dipantulkan ulang). \qty{1.2}{\micro s}:
\qty{120}{m}; gema yang tegak: sebuah rangkaian terbuka --- kabel yang terputus.
\end{solution}

\begin{solution}{exo:b2:wave-interfaces:10}
(a) $\vect E = E\vect e_z$; $\vect k\wedge\vect e_z = (k_y, -k_x, 0)$, jadi komponen tangensialnya
$B_y = -(n/c)\cos\theta\,E$ bagi gelombang datang dan terusannya dan
$+(n_1/c)\cos\theta_1E_r$ bagi yang terpantul. (b) $1 + r = t$ dan $n_1\cos\theta_1
(1 - r) = n_2\cos\theta_2t$: $r_s$ seperti yang dinyatakan. (c) $\theta = 0$: $(n_1 - n_2)/(n_1 +
n_2)$; $\theta_1 \to 90^\circ$: $r_s \to -1$. (d) $\cos\theta_2 = 0.882$: $r_s = -0.30$, $R_s =
9.2\%$; $R_p = 0.85\%$; melawan $4.3\%$ pada datang tegak lurus.
\end{solution}

\begin{solution}{exo:b2:wave-interfaces:11}
(a) \qty{270}{nm}. (b) $d = 0.35\ell = \qty{95}{nm}$; $2.3\ell = \qty{620}{nm}$. (c) Punggung
sidiknya menyentuh kacanya lalu meloloskan cahayanya ke dalam kulitnya: punggung
yang gelap, lembah yang terang pada pengindranya. (d) \omterm{def:b2:dispersion-wave-packets:complex}{Gelombang evanesennya} adalah
fungsi gelombang yang meluruh eksponensial dalam daerah terlarangnya; medium
kedua dalam panjang peluruhannya memungutnya --- penerowongan optik.
\end{solution}

\begin{solution}{exo:b2:wave-interfaces:12}
(a) Simpul $\vect E$ pada kedua cerminnya: $\sin kL = 0$, $L = m\lambda/2$, $f_m = mc/2L$.
(b) $1$, $2$, \qty{3}{GHz}: tak ada yang pada \qty{2.45}{GHz} --- ragam tiga dimensi
sebuah oven yang sungguhan jauh lebih rapat. (c) $j_s = H_0 = E_0/\mu_0c = \qty{53}{A/m}$
pada masing-masingnya, dan tekanannya $\mu_0j_s^2/2 = \qty{1.8}{mPa}$ mendorong cerminnya
saling menjauh. (d) $\delta = \qty{1.3}{\micro m}$, $R_s = \qty{12.5}{m\ohm}$; $\tfrac12R_sH_0^2 =
\qty{18}{W/m^2}$ tiap cermin; energi tersimpan reratanya $\tfrac14\varepsilon_0E_0^2 \times L = \qty{1.3e-7}{J/m^2}$;
$Q = 2\pi f \times 1.3 \times 10^{-7}/35 \approx 6 \times 10^4$.
\end{solution}

\begin{solution}{pb:b2:wave-interfaces:1}
\textbf{1.} $R = 4.3\%$; $0.957^{16} = 0.50$; separuh cahayanya tersesat.

\textbf{2.} $n = 1.23$, $e = \qty{112}{nm}$; tak ada zat padat yang awet berindeks
serendah itu (MgF$_2$, $1.38$, yang terendah): sisanya $1.3\%$.

\textbf{3.} $0.987^{16} = 0.81$; $0.997^{16} = 0.95$. Pantulan sesatnya mengaburkan
citranya: silau selubung, kontras rendah.

\textbf{4.} $\varphi = \pi \times 550/440 = 1.25\pi$: $R \approx 0.0256 + 0.0025 + 2 \times 0.008 \times
\cos1.25\pi = 1.7\%$ --- lebih besar daripada pada \qty{550}{nm}, seperti di ujung merahnya: sebuah
pantulan magenta.

\textbf{5.} Tiap antarmukanya memantulkan $|r| = 0.97/3.73 = 0.26$; jarak
seperempat gelombangnya menaruh pantulan berturut-turutnya sefase, sehingga amplitudonya bertambah
alih-alih meniadakan; reflektansi tumpukannya adalah $1 - 4(n_L/n_H)^{2N}
\approx 1 - 10^{-9}$ bagi $N = 20$.

\textbf{6.} Syarat sefasenya hanya berlaku dekat panjang gelombang
rancangannya (sebuah pita sekitar $\pm15\%$); sisanya lolos: sebuah cermin hijau
tampak magenta dalam penerusan.

\textbf{7.} Ke dalam gelombang terusannya: interferensinya menata ulang,
$R + T = 1$.

\textbf{8.} $19\%$; dengan $Z_m = \sqrt{Z_1Z_2} = \qty{6.9e6}{kg.m^{-2}.s^{-1}}$: $100\%$ pada
\qty{5}{MHz}.

\textbf{9.} $r_1 = 0.63$, $r_2 = 0.62$; pada \qty{3}{MHz} fase pulang perginya
$0.6\pi$, pada \qty{7}{MHz} $1.4\pi$: $R \approx 0.39 + 0.39 - 0.24 = 0.54$, $T \approx 0.46$
pada keduanya: penyesuaiannya hanya berlaku pada pita di sekitar \qty{5}{MHz}, sementara
pulsa pendek menuntut pita lebar --- sebuah kompromi.

\textbf{10.} $t_1 = 2 \times 410/30 \times 10^6 = 2.7 \times 10^{-5}$, $t_2 = 2.0$: $t = 5.5 \times 10^{-5}$,
$T = t^2Z_1/Z_3 = 5.6 \times 10^{-8}$: \qty{-72}{dB}.

\textbf{11.} $t = 0.095 \times 1.03 = 0.098$, $T = 0.098^2 \times 30/1.6 = 0.18$: gelnya
sekadar menyingkirkan udaranya.

\textbf{12.} $0.9^N = 0.01$: $N = 44$ perjalanan pulang pergi --- dering turun yang panjang;
penyangganya menyerap gelombang mundurnya lalu mematikannya.

\textbf{13.} $R = 0.40$ di tulangnya; $2 \times 5 \times 0.5 \times 5 = \qty{25}{dB}$
\omterm{def:b2:dispersion-wave-packets:complex}{pelemahan}: $0.40 \times 3.2 \times 10^{-3} = 1.3 \times 10^{-3}$ (\qty{-29}{dB}); tulangnya
menyerap apa yang tak dipantulkannya: tak ada yang kembali dari belakangnya.

\textbf{14.} $R = 1$, $r = -1$.

\textbf{15.} Hubung singkat di $x = 0$: $v = -2\iu v_i\sin kx$, $i = (2v_i/Z_0)\cos kx$;
di $x = -d$: $v/i = \iu Z_0\tan kd$; tak hingga pada $d = \lambda/4$.

\textbf{16.} $R_s \parallel \infty = R_s$: $r = (R_s - Z_0)/(R_s + Z_0) = 0$ bagi
\qty{377}{\ohm}; $d = \qty{7.5}{mm}$.

\textbf{17.} Lembarnya melihat medan datang penuhnya: $|v|^2/2R_s =
E_0^2/2Z_0$, yaitu intensitas datangnya --- semuanya terserap.

\textbf{18.} \qty{8}{GHz}: $kd = 72^\circ$, $Z_{\text{msk}} = \iu\qty{1160}{\ohm}$, $R_s \parallel
Z_{\text{msk}} = (340 + 111\iu)\,\Omega$, $R = |r|^2 = 2.6\%$ (\qty{-16}{dB}); \qty{12}{GHz}:
$kd = 108^\circ$, modulus yang sama: \qty{-16}{dB}. Berguna sekitar $\pm20\%$.

\textbf{19.} Datang miring memendekkan seperempat gelombang efektifnya
($d\cos\theta$) lalu mengubah impedansi gelombangnya dengan sudutnya:
penyesuaiannya hilang.

\textbf{20.} Satu tabir bekerja pada satu frekuensi dan satu sudut; lapisan
berugi yang bertingkat dan permukaan yang berbentuk khusus menyerap dan membelokkan pada pita lebar.

\textbf{21.} $\underline n = -\iu\sqrt{(\omega_p/\omega)^2 - 1} = -3.6\iu$; $|r| = 1$, fasenya
$2\arctan3.6 = \qty{149}{\degree}$; kedalamannya $c/3.6\omega = \qty{22}{nm}$. Gelombangnya tak pernah
memasuki logamnya untuk dibelah menjadi dua pantulan yang sebanding: tak ada
apa pun yang dapat ditiadakan \omterm{prop:b2:wave-interfaces:optics}{lapisan seperempat gelombang} terhadap pantulan bermodulus satu.

\textbf{22.} \qty{30}{W}; $30/(0.01 \times 235) = \qty{13}{K/s}$.

\textbf{23.} Simpulnya tiap $\lambda/2 = \qty{250}{nm}$; $E \propto \cos(kx + \varphi/2)$ dengan $\varphi =
149^\circ$: simpul pertamanya di $x = -(90^\circ + 74.5^\circ)\lambda/360^\circ = \qty{-230}{nm}$; sebutir
debu di sana duduk dalam gelap.

\textbf{24.} $I = \qty{1e7}{W/m^2}$, $E_0 = \sqrt{2I/\varepsilon_0c} = \qty{87}{kV/m}$, $j_s =
2E_0/\mu_0c = \qty{460}{A/m}$.

\textbf{25.} Logam: \omterm{prop:b2:maxwell-equations:boundary}{arus permukaan} dalam satu \omterm{def:b2:dispersion-wave-packets:complex}{kedalaman kulit}, $r \approx -1$, yang dibatasi
rugi Joule arus itu. \omterm{def:b2:waves-in-media:plasma}{Plasma}: penembusan evanesen,
$|r| = 1$ dengan fase yang bergantung frekuensi, yang dibatasi tumbukan.
Tumpukan dielektrik: banyak pantulan lemah yang bertambah sefase, $r \to 1$ dalam sebuah
pita, yang dibatasi jumlah lapisan dan lebar pitanya.
\end{solution}
